% =====================================================================
%  future.tex ------ 卷末：今後の方向と課題（今后的方向与课题）
%  原书 pp.127--133（PDF 147--153）；p.134（PDF 154）为空白页，不设页标记
% =====================================================================

\chapter*{今后的方向与课题}
\addcontentsline{toc}{chapter}{今后的方向与课题}
\markboth{今后的方向与课题}{今后的方向与课题}

% [原文 p.127]
\para{共形场理论的发展}

关于共形场理论，本书主要解说了亏格 $0$ 的情形；§1.4 所述共形块空间的定义，则由土屋--上野--山田 [45] 对一般亏格的 Riemann 面作了表述。由此，共形场理论便可看作 Riemann 面模空间上向量丛的理论。而且，这个向量丛上已构造出射影联络。这基本上是按下述方式得到的：在模空间的边界处引入 KZ 联络，再把它延拓到整个模空间。另一方面，如 §2.5 所述，共形块空间也可表述为 Riemann 面上 $G$ 平坦丛的模空间上复线丛的全纯截面空间。[17] 从 Chern--Simons 规范理论的几何量子化的立场出发，构造了以共形块空间为纤维的、Riemann 面模空间上向量丛的射影联络。

本书用共形场理论表述了三维流形的 Witten 不变量 $Z_k(M)$。由共形场理论的立场说明 Chern--Simons 理论所暗示的 $Z_k(M)$ 的渐近行为，是一个重要问题。最近，吉田朋好 [47] 实现了 Hitchin 等人尝试过的共形场理论 Abel 化纲领。其内容是：提升到 Riemann 面的二重分支覆盖上，利用 Prym 簇的理论，用 theta 函数的语言完全描述共形块空间。由此，关于射影联络的和乐结构，可以给出用 theta 函数的明确描述。

亏格 $0$ 时，KZ 联络的单值表示经 [34]、[44] 等，一般由量子群描述，这一点已由 Drinfel'd [24] 证明。此后，共形场理论与绝对 Galois 群 $\mathrm{Gal}(\overline{\mathbb{Q}}/\mathbb{Q})$ 的关系得到了深入研究。

% [Tutor Note: 所谓 Abel 化，是把原本要靠非 Abel 对象（共形块、射影联络的和乐）描述的量，搬到含阿贝尔簇（Prym 簇）的 theta 函数语言中去；二重分支覆盖正是把模空间问题化为 Prym 簇上 theta 函数问题的桥梁。]

% [原文 p.128]
\para{不变量的几何侧面与组合侧面}

本书在构造 Witten 不变量时，强调了它与共形场理论的关系等几何侧面；但仅就不变量的定义而言，组合的途径也是可能的。level $k$ 为奇数时，二阶循环群 $\Z/2\Z$ 通过 $\mu\mapsto k-\mu$ 无不动点地作用在权集合 $P_{+}(k)$ 上，据此 $Z_k(M)$ 分解为由 $M$ 的同调等拓扑量决定的项与本质项的乘积。把这个本质部分记作 $\tau_r(M)$。这里令 $r=k+2$。取 $M$ 为同调 3 球面。此时若 $r$ 为奇素数，则可证明 $\tau_r(M)\in\Z[q]$，其中 $q=e^{2\pi i/r}$。于是把 $\tau_r(M)$ 按 $q-1$ 的幂展开，再考虑 $r$ 充分大的情形，据大槻 [42]，便构造出如下取值于形式幂级数的 $M$ 的不变量：
\[
  \tau(M)=\sum_{n=0}^{\infty}\lambda_n(M)(t-1)^{n}.
\]
据村上的结果 [41]，还可知：
\[
  \lambda_0(M)=1,\qquad \lambda_1(M)=6\lambda(M).
\]
这里 $\lambda(M)$ 是 $M$ 的 Casson 不变量。这样的数论式展开，与 $Z_k(M)$ 在 $k\to\infty$ 时的渐近展开之间的关系，目前似乎尚未弄清。这是今后的重要课题。

% [Tutor Note: 记号 $\tau$ 在本书中指向多种互不相同的对象：p.128 的 $\tau_r(M)$、$\tau(M)$ 是大槻型不变量，p.131 的 $\tau$ 是 $A(\Sigma)\to C^{\infty}(M_{\Sigma})$ 的映射，p.133 的 $\tau$ 函数则属于 KdV 方程；须按上下文区分，勿当成同一记号。]

又，[39] 中对给定 framing 的链环，用 Drinfel'd 结合子构造取值于弦图的 Vassiliev 不变量，由此构造了对该链环作 Dehn 手术后所得三维流形的拓扑不变量。它以图 3.5 所示关系式为模、取值于 Feynman 图的空间，在所谓有限型不变量之中是万有的。包括它与上面的 $\tau(M)$ 的关系等在内，已有系统的研究。由 Chern--Simons 微扰理论定义的不变量，
% [原文 p.129]
取值于同一空间，但后者由 $M$ 上的平坦联络决定。两者的确切关系目前似乎尚不清楚。

纽结的 Jones 多项式的发现，历史上源自算子环。Jones 在研究 $\mathrm{II}_{1}$ 子因子 $N\subset M$ 的指数的过程中，发现了辫群的表示及其上的 Markov 迹，这成为定义 Jones 多项式的契机。$\mathrm{II}_{1}$ 子因子的分类理论此后由 Ocneanu 等人精细化，它与称为 paragroup 的组合对象之间的对应也变得清楚。其中出现了与共形场理论类似的融合代数、$6j$ 符号等结构。$6j$ 符号在共形场理论中，以 §2.1 图 2.2 所示连接矩阵的形式登场。另一方面，Turaev 与 Viro 发现了如下手法：把三维流形作三角剖分，对每个四面体配以 $6j$ 符号，再取这些乘积的状态和，从而构造三维流形的拓扑不变量。对 paragroup 也可用这一手法构造三维流形的拓扑不变量。详见 Evans--Kawahigashi [25]。阐明这类拓扑不变量与 Chern--Simons 理论之间的内在关系，当是今后的课题。

\para{用积分表示的拓扑不变量}

由 Chern--Simons 微扰理论，可以得到种种拓扑不变量的积分表示。这似将为积分几何带来新的视角。作为鲜明的例子，这里举出 $\R^{3}$ 内纽结的一个不变量：它由 Guadagnini--Martellini--Mintchev [29] 从 Chern--Simons 泛函配分函数的微扰计算得到。把 $\R^{3}\setminus\{0\}$ 上的 $\omega$ 取作式 (3.26) 所定义的 Green 形式。$\omega$ 是 $\R^{3}\setminus\{0\}$ 上的 2 次微分形式，满足 $\int_{S^{2}}\omega=1$，并定出球面 $S^{2}$ 的体积元。对由嵌入 $f\colon S^{1}\to\R^{3}$ 表示的纽结 $K$，把 $S^{1}$ 参数表示为 $e^{2\pi it}$，$0\le t\le 1$，考虑如下积分：
\[
  I_X=\int_{0<t_{1}<t_{2}<t_{3}<t_{4}<1}
  \omega\bigl(f(t_{1})-f(t_{3})\bigr)\wedge\omega\bigl(f(t_{2})-f(t_{4})\bigr),
\]
% [原文 p.130]
\[
  I_Y=\int_{0<t_{1}<t_{2}<t_{3}<1,\ x\in\R^{3}\setminus K}
  \omega\bigl(f(t_{1})-x\bigr)\wedge\omega\bigl(f(t_{2})-x\bigr)\wedge\omega\bigl(f(t_{3})-x\bigr).
\]
$I_X$、$I_Y$ 分别是纽结全体空间上的函数，而
\[
  I_X+I_Y-\frac{1}{24}
\]
是取整数值的纽结拓扑不变量，可以证明它与 §3.1 所述 Conway 多项式的 2 次项系数 $a_2(K)$ 一致。在纽结全体的空间上考虑以适当的点构形空间为纤维的纤维丛，则 $I_X$、$I_Y$ 可分别解释为把 Green 形式的积拉回后在纤维上积分所得。要证明拓扑不变性，只需在纽结全体的空间上证明 $d(I_X+I_Y)=0$；而左端外微分的项，可通过点构形空间的紧化来描述。这里，对图作使边退化的操作以定义边界算子所得的图复形起重要作用。详见 Bott--Taubes [21] 等。

% [Tutor Note: 式中的 $1/24$ 是 Chern--Simons 微扰计算给出的常数项（framing 修正的一部分），与 §1.1 式 (1.8) 的 $1/24\pi^{2}$、式 (1.7) 的 $1/4\pi^{2}$ 没有直接关系；不可因数字相同而把它们对应起来。]

在积分 $I_X$、$I_Y$ 中，让纽结逼近其投影图，并把体积元 $\omega$ 集中到球面的两极，则 $I_X$ 一方成为仅由纽结交叉点决定的量。另一方面，$I_Y$ 一方趋于由纽结投影图中忽略交叉的上下所得到的平面曲线决定的量。这个量在按奇异性把平面曲线全体空间分割时，于余维数 $1$ 的成分上取常值。这样的平面曲线不变量，已由 Arnold [13] 从接触几何的立场作了系统研究。由 Chern--Simons 微扰理论得到的、用积分表示的拓扑不变量，与 Arnold 的平面曲线不变量之间的关系，是有趣的问题。

第 3 章中，Vassiliev 不变量的记述用了弦图。从 Chern--Simons 微扰理论的立场看，这对应于不变量的渐近展开中出现的 Feynman 图。取 $\Sigma$ 为 Riemann 面，则 §3.1 的构造可推广为 $\Sigma$ 上弦图的 Vassiliev 不变量与 $\Sigma\times I$ 内链环的 Vassiliev 不变量。把 $\Sigma$ 上弦图的同伦类全体张成的复线性空间按 4 项关系式作商，记作 $A(\Sigma)$，
% [原文 p.131]
则 $A(\Sigma)$ 上自然可以引入 Poisson 代数结构。这是把对 $\Sigma$ 上环路的同伦类定义的 Goldman 括号自然推广所得。对 Lie 群 $G=SU(2)$，把 $\Sigma$ 上平坦 $G$ 丛的模空间记作 $M_{\Sigma}$。$M_{\Sigma}$ 具有辛结构。把 $G$ 的表示对应于弦图的环路，并在 §3.1 的权系定义中加上沿环路的 $G$ 丛的和乐信息，便得到从 $A(\Sigma)$ 到 $M_{\Sigma}$ 上函数空间 $C^{\infty}(M_{\Sigma})$ 的映射 $\tau$。可以证明它是 Poisson 代数的同态。从 $\Sigma\times I$ 内具有二重点的链环全体的线性空间到 $A(\Sigma)$ 有投影，把上述 Poisson 代数同态提升，便可构造取值于 $C^{\infty}(M_{\Sigma})$ 的某个非交换化的、$\Sigma\times I$ 内链环的 Vassiliev 不变量。

% [Tutor Note: 4 项关系式即弦图的 4T 关系；按它作商之后 $A(\Sigma)$ 上才可能有 Poisson 结构，这与第 3 章用弦图定义 Vassiliev 不变量时所用的关系式是同一套。]

共形块空间经过几何量子化的过程，可以看作 $C^{\infty}(M_{\Sigma})$ 的非交换化。把 Chern--Simons 微扰理论对带边界的三维流形展开时，平坦联络的处理成为问题，而上述手法或许给出一种途径。又，若取 $G=SL(2,\R)$，则上述 $\tau$ 对环路所给出的是：在相应的 $\Sigma$ 的双曲几何结构中，把环路在同伦范围内实现为最短测地线时，该测地线的长度。本书处理的是紧群的 Chern--Simons 理论；而展开非紧群的理论，对讨论与双曲几何的关联是重要的。

\para{与曲面模空间的关系}

先说明矩阵积分。设 $H_N$ 为 $N$ 阶 Hermite 矩阵 $M=(M_{ij})$ 全体。对矩阵的多项式 $G(X)$，考虑
\[
  \int_{H_N}\exp\bigl(-\Tr G(M)\bigr)\,dM
\]
这种形状的 Hermite 矩阵积分。这里 $dM$ 是 $N^{2}$ 维 Euclid 空间 $H_N$ 的 Euclid 不变 Lebesgue 测度，
% [原文 p.132]
并作归一化使
\[
  \int_{H_N}\exp\bigl(-\Tr(M^{2}/2)\bigr)\,dM=1
\]
成立。此处值得注意的是矩阵积分与带状图个数的母函数之间的关系。这里，带状图指对汇聚于各顶点的边给定循环顺序的图。

作为例子，取
\[
  G(M)=\frac{M^{2}}{2}+\frac{tM^{4}}{4N},
\]
此时把上述矩阵积分记作 $Z(N,t)$，试计算关于 $t$ 的展开。一般地令
\[
  \bigl\langle M_{i_{1}j_{1}}\cdots M_{i_{n}j_{n}}\bigr\rangle
  =\int_{H_N}M_{i_{1}j_{1}}\cdots M_{i_{n}j_{n}}
  \exp\bigl(-\Tr(M^{2}/2)\bigr)\,dM,
\]
按与 §3.3 相同的方法计算，得 $\langle M_{ij}M_{kl}\rangle=\delta_{il}\delta_{jk}$。因此，$\bigl(\Tr(M^{4})\bigr)^{m}$ 的展开中出现的每一项，对应于顶点 $m$ 个、每个顶点周围汇集 4 条边的带状图。在带状图的边界上粘合圆盘，便可构成 Riemann 面。设 $W(g,m)$ 为上述带状图中给出亏格 $g$ 的 Riemann 面的个数，并考虑图计数中的重复，则得
\[
  \log Z(N,t)=\sum_{g=0}^{\infty}N^{2-2g}\sum_{m=0}^{\infty}(-1)^{m}W(g,m)t^{m}.
\]
考虑曲面上带状图的对偶剖分，则 $W(g,m)$ 也可看作把亏格 $g$ 的闭曲面用 $m$ 个四边形剖分的个数。

% [Tutor Note: 展开式中 $N^{2-2g}$ 的指数就是亏格 $g$ 闭曲面的 Euler 示性数：每增加一个把手，$N$ 的幂次降 2。这正是把矩阵积分读成曲面计数（大 $N$ 展开）的机制。]

把亏格 $g$、带 $n$ 个点的 Riemann 面的模空间记作 $M_{g,n}$，则带状图对应于 $M_{g,n}\times\R_{+}^{n}$ 的胞腔分解中的胞腔。这里 $n$ 对应于粘合到带状图上的圆盘个数。这个胞腔分解是 Mumford、Harer 等人用二次微分（quadratic differential）的理论作出的。给带状图的边赋以长度信息，就构造出 Riemann 面的复结构。

% [Tutor Note: 这里的「二次微分」是 quadratic differential，即 $dz^{2}$ 型的张量，与 2 次微分形式（2-form）不是同一对象；Mumford--Harer 的胞腔分解正是用它（Strebel 微分）的水平轨线作出的。]

% [原文 p.133]
1990 年前后，Witten [46] 针对二维量子重力理论的 Lagrangian，提出了计算 $n$ 点函数的两种方法。第一种方法是上面说明过的、用矩阵积分求 Riemann 面上带状图个数的方法；第二种方法是利用 Riemann 面模空间的相交理论的方法。由此表述了一个猜想：两种方法所得的 $n$ 点函数一致。如 [40] 所解说，矩阵积分与可积系统的理论深刻相关。Witten 的这个猜想也暗示了模空间的相交理论与 KdV 方程的 $\tau$ 函数之间的关系。Witten 的猜想后来由 Kontsevich [37] 解决。另一方面，把目光转向 Chern--Simons 微扰理论，便会注意到：构造拓扑不变量所用的 2 圈 Feynman 图（图 3.4），无非就是 $M_{0,3}\times\R_{+}^{3}$ 的胞腔分解中出现的带状图。证明 Chern--Simons 微扰理论中出现的积分不依赖度量，基本上要用与胞腔分解相对应的图的相邻关系。这些预示着一股卷入拓扑场论、可积系统、Morse 同伦理论等在内的更大潮流。

% CHUNK future DONE
